\documentclass[11pt]{amsart}
\usepackage[T1]{fontenc}
\usepackage{lmodern}
\usepackage{amsmath,amssymb,amsthm}
\usepackage[margin=1.1in]{geometry}
\usepackage{booktabs}
\usepackage{hyperref}

\newtheorem{theorem}{Theorem}[section]
\newtheorem{lemma}[theorem]{Lemma}
\newtheorem{proposition}[theorem]{Proposition}
\newtheorem{corollary}[theorem]{Corollary}
\newtheorem{conjecture}[theorem]{Conjecture}
\newtheorem{question}[theorem]{Question}
\theoremstyle{definition}
\newtheorem{definition}[theorem]{Definition}
\newtheorem{remark}[theorem]{Remark}
\newtheorem{example}[theorem]{Example}

\newcommand{\cD}{\mathcal D}

\title[Prescribed end edges]{Prescribed end edges in Lov\'asz decompositions\\ with one even vertex}
\author{Justin Williams}
\date{November 2026 (working draft)}

\begin{document}
\begin{abstract}
Let $K$ be a graph with at most one vertex $z$ of positive even degree. A \emph{Lov\'asz
decomposition} of $K$ is a partition of $E(K)$ into paths in which every vertex $v$ is the end of
exactly $d(v)\bmod 2$ paths; Lov\'asz proved that one exists. For an edge $xy$ with $d(x)$ odd, say
that $xy$ is \emph{terminal at $x$} if the path ending at $x$ has $xy$ as its last edge. When all
degrees are odd, every edge is terminal at either end in some Lov\'asz decomposition. We ask
whether this survives one even vertex.

\begin{itemize}
\item \emph{Away from $z$ it does not.} On five vertices there is an edge $xy$, $y\ne z$, that is
terminal at $x$ in no Lov\'asz decomposition. Computation finds $1{,}410$ such edges among the
$330{,}449$ edges $xy$ ($y\ne z$, $d(x)\ge3$) of the $11{,}936$ connected graphs of order~$9$ with one
even vertex.
\item \emph{At $z$ we conjecture that it does} (Conjecture~EE: every edge $uz$ is terminal at $u$),
with no counterexample among all graphs of order at most $9$, more than $30$ million edges of graphs of
order $11$, and random graphs of orders $11$--$17$. We prove EE whenever $u$ and $z$ have at most one
common neighbour, for every degree of $u$. This contains all earlier partial results (no triangle on
$uz$; $d(u)=3$ outside two configurations) and the whole case $d(z)=2$.
\item \emph{A protected-end lemma.} Fan's transfer lemma never moves the end of a vertex $x$ not
adjacent to the target, nor of a vertex whose end edge goes to the target. With it, Lov\'asz's
Case~I step reduces a prescription away from $z$ to EE at another vertex. Hence: if some $z$--$y$
path has all interior vertices outside $N[x]$, then EE implies that $xy$ is terminal at $x$.
Unconditionally, $xy$ is terminal at $x$ when $y\in N(z)$ and $x,y$ have at most one common
neighbour other than $z$. Every one of the failures found by computation lies outside the reach of
this reduction.
\item \emph{The remaining case of EE} ($u,z$ with two or more common neighbours) reduces, through a
triangle $uzw$, to an avoidance statement of $L_1$ type, which we isolate as a conjecture ($L_1^*$)
together with its computational evidence.
\end{itemize}
The main results are formalised in Lean~4 with Mathlib, with no \texttt{sorry} and only the standard
axioms. Computations are in a separate, labelled section; no proof depends on them.
\end{abstract}
\maketitle

\section{Introduction}\label{sec:intro}

All graphs are finite and simple. A \emph{path decomposition} of $K$ partitions $E(K)$ into paths,
each with at least one edge. For a decomposition $\cD$ and a vertex $v$, $b_\cD(v)$ is the number of
paths of $\cD$ with an end at $v$; always $b_\cD(v)\equiv d(v)\pmod 2$. A vertex has \emph{positive
even degree} if $d(v)>0$ is even. Following~\cite{RD} we call $\cD$ a \emph{Lov\'asz decomposition}
if $b_\cD(v)=d(v)\bmod 2$ for every $v$: every odd vertex ends exactly one path and every even vertex
ends none. For an odd vertex $x$ we write $Q_x$ for the path ending at $x$.

\begin{theorem}[Lov\'asz~\cite{L68}]\label{thm:lovasz}
A graph with at most one vertex of positive even degree has a Lov\'asz decomposition.
\end{theorem}

\begin{definition}
Let $xy$ be an edge with $d(x)$ odd. Then $xy$ is \emph{terminal at $x$} in a Lov\'asz
decomposition $\cD$ if $xy$ is the edge of $Q_x$ at $x$, equivalently if $xy\in E(Q_x)$.
\end{definition}

\begin{theorem}[prescribed end edge~\cite{MS}]\label{thm:endedge}
If all degrees of $K$ are odd, every edge $xy$ is terminal at $x$ in some Lov\'asz decomposition.
\end{theorem}

The question of this paper is what remains of Theorem~\ref{thm:endedge} when $K$ has one vertex $z$
of positive even degree. There are two kinds of edges $xy$ with $d(x)$ odd: those at $z$ ($y=z$), and
those away from $z$ ($y\ne z$, so $d(y)$ is odd too).

\begin{conjecture}[PEE, prescribed end edge]\label{conj:pee}
In every graph with at most one vertex of positive even degree, every edge $xy$ with $d(x)$ odd is
terminal at $x$ in some Lov\'asz decomposition.
\end{conjecture}

\begin{conjecture}[EE~\cite{MS}]\label{conj:ee}
Let $z$ be the only vertex of positive even degree of $K$, and let $u\in N(z)$. Then $uz$ is terminal
at $u$ in some Lov\'asz decomposition of $K$.
\end{conjecture}

Conjecture~\ref{conj:ee} is the case $y=z$ of Conjecture~\ref{conj:pee}. It is the step of
Lov\'asz's induction at the even vertex with a trivial Fan chain~\cite[Proposition~EE]{MS}, and it is
the one-transition case of the star-terminal programme for the spine tests of Gallai's
conjecture~\cite{STL}. The multi-pair separation manuscript~\cite{MS} proves it when $uz$ is in no
triangle, and when $d(u)=3$ outside two local configurations; the star-terminal
manuscript~\cite{STL} adds the case $d(z)=2$ with $u$ and the other neighbour of $z$ having no second
common neighbour.

\subsection*{Results}
\begin{enumerate}
\item Conjecture~\ref{conj:pee} is false away from the even vertex (Proposition~\ref{prop:counter}). The
smallest counterexample has five vertices.
\item Conjecture~\ref{conj:ee} holds whenever $u$ and $z$ have at most one common neighbour
(Theorem~\ref{thm:common}). In particular it holds whenever $d(z)=2$ (Corollary~\ref{cor:deg2}),
which completes the star-terminal conjecture at $d(z)=2$.
\item A \emph{protected-end lemma} (Lemma~\ref{lem:protect}): Fan's transfer with target $g$ keeps a
terminal edge $xy$ at $x$ when $x\notin N(g)$ or $y=g$. Consequences
(Theorem~\ref{thm:path}, Corollary~\ref{cor:adj}): EE reaches every edge $xy$ away from $z$ for which
some $z$--$y$ path has its interior outside $N[x]$; and $xy$ is terminal at $x$ unconditionally when
$y\in N(z)$ and $x,y$ have at most one common neighbour other than $z$. If $z$ has a pendant
neighbour, Conjecture~\ref{conj:pee} holds for $K$ (Proposition~\ref{prop:pendant}).
\item The remaining case of EE reduces to $L_1$-type avoidance statements
(Theorem~\ref{thm:triangle}, Proposition~\ref{prop:moveA}), unified as Conjecture~\ref{conj:L1star}.
\item Computations (Section~\ref{sec:comp}): an exact oracle, validated against brute force and an
independent oracle; exhaustive results to order $9$ (EE to order $11$); random graphs to order~$17$.
In the data, the failures of PEE are exactly confined to the edges that the reduction of item~3
cannot reach.
\item Lean formalisation of Theorem~\ref{thm:common}, Corollary~\ref{cor:deg2}, the case $k=1$ of Corollary~\ref{cor:adj}, and Theorem~\ref{thm:triangle} (Section~\ref{sec:lean}).
\end{enumerate}

\section{Preliminaries}\label{sec:prelim}

For a decomposition $\cD$ and a vertex $g$ let $Z_\cD(g)=\{w\in N(g):b_\cD(w)=0\}$.

\begin{lemma}[{Fan's lemma with vertex sets~\cite{Fan05},~\cite[Lemma~2.2]{RD}}]\label{lem:fan}
Let $\cD$ be a path decomposition of $K$, let $t\ne g$ be non-adjacent, and suppose
$b_\cD(t)>|Z_\cD(g)|$. Then $K+tg$ has a path decomposition $\cD^+$ with
$b_{\cD^+}=b_\cD-\mathbf 1_t+\mathbf 1_g$ in which every path has its vertex set inside $V(Q)\cup\{g\}$
for some path $Q$ of $\cD$.
\end{lemma}

\begin{theorem}[{separation~\cite[Theorem~6.1]{RD}}]\label{thm:sep}
If $K$ has at most one vertex of positive even degree and $c,a$ are at distance at least three, then
some Lov\'asz decomposition of $K$ has no path containing both $c$ and $a$.
\end{theorem}

\begin{proposition}[{\cite[Proposition~EE]{MS}, no classes}]\label{prop:EEequiv}
Let $z$ be the only vertex of positive even degree, $u\in N(z)$ and $K'=K-uz$. Then $uz$ is terminal at
$u$ in some Lov\'asz decomposition of $K$ if and only if $K'$ has a Lov\'asz decomposition whose path
$Q_z$ avoids $u$. (Append $zu$ to $Q_z$; conversely delete $uz$.)
\end{proposition}

\emph{The leaf graph.} For an edge $xy$ let $H(K;x,y)$ be $K-xy$ together with a new vertex $\ell$ and
the edge $y\ell$. Identifying $\ell$ with $x$ maps $H(K;x,y)$ onto $K$, bijectively on edges.

\begin{lemma}[leaf merge~\cite{MS}]\label{lem:merge}
Let $\cD_H$ be a path decomposition of $H=H(K;x,y)$ in which no path contains both $x$ and $\ell$.
Identifying $\ell$ with $x$ gives a path decomposition $\cD$ of $K$ with $b_\cD(x)=b_{\cD_H}(x)+
b_{\cD_H}(\ell)$ and $b_\cD(w)=b_{\cD_H}(w)$ otherwise. The path through $y\ell$ becomes a path through
$yx$ with the same ends, $\ell$ replaced by $x$.
\end{lemma}

\begin{lemma}\label{lem:leafequiv}
Let $d(x)$ be odd. Then $xy$ is terminal at $x$ in some Lov\'asz decomposition of $K$ if and only if
$H(K;x,y)$ has a decomposition $\cD_H$ with $b_{\cD_H}(x)=0$, $b_{\cD_H}(\ell)=1$ and
$b_{\cD_H}(w)=d_K(w)\bmod 2$ for $w\ne x,\ell$, in which no path contains both $x$ and $\ell$. We call
such a $\cD_H$ a \emph{leaf witness} for $xy$.
\end{lemma}

\begin{proof}
Given $\cD$, write $Q_x=Ryx$; $x$ occurs on $Q_x$ only at its end. Replace $Q_x$ by $Ry\ell$. The
other paths do not contain $\ell$, and $Ry\ell$ does not contain $x$. The end counts are as stated.
Conversely apply Lemma~\ref{lem:merge}: $b_\cD(x)=0+1$, and the path through $y\ell$, which ends at
$\ell$ since $\ell$ is a leaf, becomes a path ending at $x$ through $yx$.
\end{proof}

In the leaf graph, ``$xy$ is terminal at $x$'' has become a property of vertex sets, and
Lemma~\ref{lem:fan} preserves it whenever the target is neither $x$ nor $\ell$. This is the engine of
everything below.

\section{Away from the even vertex, PEE fails}\label{sec:counter}

\begin{proposition}\label{prop:counter}
Let $z$ be the only vertex of positive even degree of $K$. Suppose $d(z)=2$ with $N(z)=\{x,t\}$ and
$d(x)=3$ with $N(x)=\{z,t,y\}$. Then $xy$ is terminal at $x$ in no Lov\'asz decomposition. Such graphs
exist on five vertices: $x,t,z,y,p$ with edges $xt,xz,xy,tz,ty,yp$.
\end{proposition}

\begin{proof}
Suppose $xy$ is terminal at $x$. Then $x$ has its only end on $Q_x$, through $xy$; the other two edges
$xz$, $xt$ at $x$ are consecutive on one path $P$. Since $z$ has no end and degree two, $zx$ and $zt$
are consecutive on a path; this is $P$, which contains $zx$. So $P$ contains $t\,x\,z\,t$, which repeats
$t$. In the example the degrees are $d(x)=d(t)=d(y)=3$, $d(z)=2$, $d(p)=1$.
\end{proof}

So Conjecture~\ref{conj:pee} fails for edges away from $z$. The computations of
Section~\ref{sec:comp} show many more failures, of several local types, with $d(x)\in\{3,5,7\}$ and
$d(z)\in\{2,4,6\}$. In every failure found, $y\notin N(z)$, $\mathrm{dist}(x,z)\le2$, and $xy$ is
terminal at $y$ in some Lov\'asz decomposition. Section~\ref{sec:away} explains the first property.
The other two remain conjectural (Question~\ref{q:away}).

\section{The protected-end lemma}\label{sec:protect}

\begin{lemma}[protected end]\label{lem:protect}
Let $\cD$ be a decomposition of $K$, let $t\ne g$ be non-adjacent, and let $x\notin\{t,g\}$ have
$b_\cD(x)=1$, with $Q_x$ ending at $x$ through the edge $xy$. Suppose that either $x\notin N(g)$ or
$y=g$, and that $b_\cD(t)>|Z_\cD(g)|$. Then $K+tg$ has a decomposition $\cD^+$ with
$b_{\cD^+}=b_\cD-\mathbf 1_t+\mathbf 1_g$ whose path ending at $x$ goes through $xy$.
\end{lemma}

\begin{proof}
Pass to the leaf graph $H=H(K;x,y)$ as in the proof of Lemma~\ref{lem:leafequiv}: replace $Q_x=Ryx$ by
$Ry\ell$. This gives a decomposition $\cD_H$ of $H$ with $b_{\cD_H}(x)=0$, $b_{\cD_H}(\ell)=1$ and the
other end counts of $\cD$, in which no path contains both $x$ and $\ell$. In $H$ the vertex $x$ is not
adjacent to $g$: either $x\notin N_K(g)$, or $y=g$ and the edge $xg$ has been removed. The leaf $\ell$
has an end. Hence $Z_{\cD_H}(g)\subseteq Z_\cD(g)$ and $b_{\cD_H}(t)=b_\cD(t)$ (as $t\ne x$), so
Lemma~\ref{lem:fan} applies in $H$ to $t$ and $g$ (non-adjacent in $H$, since $H-\ell\subseteq K$).
It gives a decomposition of $H+tg=H(K+tg;x,y)$ with $b=b_{\cD_H}-\mathbf 1_t+\mathbf 1_g$ in which every
path lies in an old path together with $g$. Since $g\notin\{x,\ell\}$, still no path contains both $x$
and $\ell$. Lemma~\ref{lem:merge} gives $\cD^+$.
\end{proof}

\begin{remark}
Lemma~\ref{lem:protect} can also be read off the chain in the proof of Fan's lemma~\cite[Lemma~2.2]{RD}.
Every chain vertex after $t$ is a vertex $\pi$ immediately before $g$ on a path read from one of its
ends. A vertex $x\notin N(g)$ is never one. Neither is a vertex whose end path has $xg$ as its end
edge: on that path $x$ comes after $g$ when it is read from its other end. A re-routing changes a path
only near the end it is read from, and the append happens at a chain vertex, so the end at $x$ and its
edge are untouched. The leaf proof uses only the statement of Lemma~\ref{lem:fan}, and it is the one
formalised.
\end{remark}

\section{EE when \texorpdfstring{$u$}{u} and \texorpdfstring{$z$}{z} have at most one common
neighbour}\label{sec:common}

\begin{theorem}\label{thm:common}
Let $z$ be the only vertex of positive even degree of $K$ and $u\in N(z)$. If $u$ and $z$ have at
most one common neighbour, then $K-uz$ has a Lov\'asz decomposition in which no path contains both
$u$ and $z$. Consequently $uz$ is terminal at $u$ in some Lov\'asz decomposition of $K$.
\end{theorem}

\begin{proof}
If $u$ and $z$ have no common neighbour they are at distance at least three in $K'=K-uz$, and
Theorem~\ref{thm:sep} applies ($K'$ has at most one vertex of positive even degree, namely $u$).

Let $w$ be the unique common neighbour, and let $K_0=K-uz-uw$. Degrees: $d_{K_0}(u)=d(u)-2$, which is
odd as $d(u)$ is odd and at least $3$; $d_{K_0}(w)=d(w)-1$, which is even and positive, as $w\ne z$
has odd degree at least two; $d_{K_0}(z)=d(z)-1$ is odd; all other degrees are as in $K$. So $w$ is the
only vertex of $K_0$ of positive even degree. In $K_0$ the vertices $u,z$ are non-adjacent, and a
common neighbour would be a common neighbour in $K$ other than $w$; so they are at distance at least
three. Theorem~\ref{thm:sep} gives a Lov\'asz decomposition $\cD_0$ of $K_0$ with no path containing
both $u$ and $z$.

Now apply Lemma~\ref{lem:fan} to $K_0$, $t=u$, target $w$. Every neighbour of $w$ in $K_0$ has odd
degree, hence an end, so $Z_{\cD_0}(w)=\emptyset$, while $b_{\cD_0}(u)=1$. The result is a
decomposition $\cD_1$ of $K_0+uw=K-uz$ with $b_{\cD_1}(u)=0$, $b_{\cD_1}(w)=1$ and the other end counts
of $\cD_0$; these are the Lov\'asz end counts of $K-uz$. Every path of $\cD_1$ lies in a path of
$\cD_0$ together with $w\notin\{u,z\}$, so no path of $\cD_1$ contains both $u$ and $z$. In particular
$Q_z$ avoids $u$, and Proposition~\ref{prop:EEequiv} finishes the proof.
\end{proof}

\begin{corollary}\label{cor:deg2}
If the even vertex $z$ has degree two, both of its edges are terminal at their other ends in suitable
Lov\'asz decompositions. In particular, the star-terminal conjecture of~\cite{STL} holds whenever
$d(z)=2$, including the configuration with two triangles $zuw$, $tuw$ left open there.
\end{corollary}

\begin{proof}
If $N(z)=\{u,w\}$, the only possible common neighbour of $u$ and $z$ is $w$. For the second
statement: by~\cite{STL}, STL at $d(z)=2$ is EE at the other neighbour.
\end{proof}

\begin{remark}\label{rem:covers}
Theorem~\ref{thm:common} contains the earlier partial results on EE without classes: $uz$ in no
triangle; $N(u)=\{z,a,b\}$ with $ab,zb\notin E$ (then the common neighbours lie in $\{a\}$); and
$N(u)=\{z,a,b\}$, $ab\in E$, with $u$ the only common neighbour of $a,b$ (then $z$ is not adjacent to
both $a$ and $b$)~\cite[Proposition~EE3]{MS}. It also contains the two degree-two cases
of~\cite{STL}. It does not depend on $d(u)$. By Theorem~\ref{thm:common} the
remaining configurations of EE are those where $u$ and $z$ have at least two common neighbours,
which forces $d(z)\ge4$.
\end{remark}

\section{Edges away from the even vertex}\label{sec:away}

\begin{theorem}\label{thm:path}
Let $z$ be the only vertex of positive even degree of $K$, and let $xy$ be an edge with $x\ne z$ and
$y\ne z$. Suppose $z=e_0,e_1,\dots,e_k=y$ ($k\ge1$) is a path in $K$ with $e_1,\dots,e_{k-1}\notin
N[x]$, and let $K_k=K-\{e_0e_1,\dots,e_{k-1}e_k\}$. Then $y$ is the only vertex of $K_k$ of positive
even degree, $x\in N_{K_k}(y)$, and if $xy$ is terminal at $x$ in some Lov\'asz decomposition of $K_k$
(Conjecture~\ref{conj:ee} for $K_k$, $y$, $x$), then $xy$ is terminal at $x$ in some Lov\'asz
decomposition of $K$.
\end{theorem}

\begin{proof}
Let $K_i=K-\{e_0e_1,\dots,e_{i-1}e_i\}$. The path visits distinct vertices, so in $K_i$ the vertex
$e_i$ has lost one edge, $e_0,\dots,e_{i-1}$ have lost one ($e_0$) or two edges, and the others none.
So $e_i$ is the only vertex of $K_i$ of positive even degree: for $0<i$, $e_i\ne z$ is odd in $K$ and has
degree at least two (two path neighbours for $i<k$; $e_{k-1}$ and $x\ne e_{k-1}$ for $i=k$). The edge
$xy$ is not on the path, and $x$ is not a path vertex: $x\ne z,y$, and $e_1,\dots,e_{k-1}\notin N[x]$.

Let $\cD_k$ be a Lov\'asz decomposition of $K_k$ in which $xy$ is terminal at $x$. For $i=k,\dots,1$
apply Lemma~\ref{lem:protect} to $K_i$ with $t=e_{i-1}$ and target $g=e_i$. In $K_i$ we have
$b(e_{i-1})=1$, and every neighbour of $e_i$ has odd degree, so $Z(e_i)=\emptyset$. Also
$x\notin\{t,g\}$. For $i<k$ the target $e_i$ is not adjacent to $x$, and for $i=k$ the target is $y$. So
the lemma applies and gives a Lov\'asz decomposition of $K_{i-1}$ in which $xy$ is terminal at $x$.
After the step $i=1$ we reach $K_0=K$.
\end{proof}

Call an edge $xy$ away from $z$ \emph{reachable} if a path as in Theorem~\ref{thm:path} exists, that
is, if $y\in N(z)$ or $y$ has a neighbour in the component of $z$ in $K-(N[x]\setminus\{z\})$.

\begin{corollary}\label{cor:adj}
Let $z$ be the only vertex of positive even degree, let $y\in N(z)$, and let $xy$ be an edge with
$x\ne z$. If $x$ and $y$ have at most one common neighbour other than $z$, then $xy$ is terminal at
$x$ in some Lov\'asz decomposition. More generally, $xy$ is terminal at $x$ whenever it is reachable
through a path with $k\ge2$ and $x,y$ have at most one common neighbour.
\end{corollary}

\begin{proof}
Apply Theorem~\ref{thm:path} with $k=1$. In $K_1=K-zy$ the common neighbours of $x$ and $y$ are those of
$K$ other than $z$, so Theorem~\ref{thm:common} gives the hypothesis for $K_1$, $y$, $x$. For $k\ge2$ we
have $N_{K_k}(x)=N(x)$ and $N_{K_k}(y)=N(y)\setminus\{e_{k-1}\}$ with $e_{k-1}\notin N(x)$.
\end{proof}

\begin{corollary}\label{cor:EEimplies}
If Conjecture~\ref{conj:ee} holds, then every reachable edge $xy$ away from $z$ is terminal at $x$ in
some Lov\'asz decomposition. In particular this holds for every edge $xy$ with $y\in N(z)$.
\end{corollary}

Proposition~\ref{prop:counter} is an unreachable edge: the only routes from $z$ to $y$ pass through
$x$ or $t\in N(x)$. By computation (Section~\ref{sec:comp}), all $1{,}410$ failures at order $9$ and all
$56{,}739$ failures found at order $11$ are unreachable; no reachable edge failed. Not every
unreachable edge fails: at order~$9$, $1{,}410$ of the $59{,}574$ unreachable edges do.

\begin{proposition}\label{prop:pendant}
If the even vertex $z$ has a neighbour $v$ of degree one, then every edge $xy$ with $d(x)$ odd is
terminal at $x$ in some Lov\'asz decomposition of $K$.
\end{proposition}

\begin{proof}
If $x=v$ there is nothing to prove: the path ending at $v$ contains its only edge. Otherwise $K-zv$
has all degrees odd or zero, and Theorem~\ref{thm:endedge} gives a Lov\'asz decomposition of $K-zv$
in which $xy$ is terminal at $x$. Append $zv$ to the path ending at $z$; this is a path because $v$ is
isolated in $K-zv$. Now $z$ has no end and $v$ has one. The path ending at $x$ is unchanged, or, if
it was the path ending at $z$, it was extended at its other end $z\ne x$.
\end{proof}

\begin{question}\label{q:away}
Is every edge $xy$ away from $z$ terminal at $x$ or at $y$? Is every edge $xy$ with
$\mathrm{dist}(x,z)\ge3$ terminal at $x$? Both hold in all our data. The first is trivial when $d(y)=1$.
\end{question}

\section{The remaining case of EE}\label{sec:rest}

By Theorem~\ref{thm:common}, EE is open only when $u$ and $z$ have at least two common neighbours. Two
reductions cover this case. Both follow the pattern of Theorem~\ref{thm:common}: delete edges at $u$,
find a decomposition with an avoidance property, restore by Fan's lemma.

\begin{proposition}[an outside neighbour]\label{prop:moveA}
Let $z$ be the only vertex of positive even degree, $u\in N(z)$ and $w\in N(u)\setminus N[z]$. If
$K_0=K-uz-uw$ has a Lov\'asz decomposition whose path ending at $z$ avoids $u$, then $uz$ is terminal
at $u$ in some Lov\'asz decomposition of $K$.
\end{proposition}

\begin{proof}
In $K_0$ the vertex $w$ is the only one of positive even degree (or there is none), $u$ and $z$ are odd.
Attach a leaf $\ell$ at $z$ and extend $Q_z$ by $z\ell$; in the resulting decomposition of $K_0+z\ell$
no path contains both $\ell$ and $u$, and $b(z)=0$. Apply Lemma~\ref{lem:fan} with $t=u$ and target $w$:
in $K_0+z\ell$ the vertex $z$, which now has no end, is not adjacent to $w$, and every other neighbour of
$w$ is odd; so $Z(w)=\emptyset<1=b(u)$. The separation of $\ell$ and $u$ survives ($w\ne\ell,u$).
Delete $\ell$: the path through $\ell$ becomes the path ending at $z$ in a Lov\'asz decomposition of
$K-uz$, and it avoids $u$. Proposition~\ref{prop:EEequiv} finishes the proof.
\end{proof}

\begin{theorem}[triangle reduction]\label{thm:triangle}
Let $z$ be the only vertex of positive even degree, $u\in N(z)$, and $w\in N(u)\cap N(z)$. If
$K_1=K-\{uz,uw,wz\}$ has a Lov\'asz decomposition whose path ending at $w$ avoids $u$ and $z$, then
$uz$ is terminal at $u$ in some Lov\'asz decomposition of $K$.
\end{theorem}

\begin{proof}
In $K_1$ the degrees of $u,z,w$ drop by two, so the parities are those of $K$, and $w$, $u$ are odd.
Let $\cD_1$ be as in the hypothesis. Let $F=K-uz$ and $H=H(F;z,w)$: here the leaf $\ell$ hangs at $w$
and stands for $z$. Extending $Q_w$ by $w\ell$ gives a decomposition $\cD_3$ of $H-uw=K_1+w\ell$ in
which no path contains $\ell$ together with $u$ or with $z$. Its end counts are those of $\cD_1$,
except that $w$ has none and $\ell$ has one. In $H-uw$ every neighbour of $w$ other than $\ell$ has odd
degree; $u$ and $z$ are not neighbours. So Lemma~\ref{lem:fan} with $t=u$ and target $w$ restores $uw$.
The result, $\cD_H$, decomposes $H$, keeps both separations, has $b(u)=0$ and $b(w)=1$, and has the
other end counts of $\cD_3$. By Lemma~\ref{lem:merge} (with $x=z$), $\cD_H$ gives a decomposition $\cD$
of $F$ with $b_\cD(z)=0+1$, which is a Lov\'asz decomposition of $F$. Its path ending at $z$ is the
image of the path through $\ell$; the latter avoids $u$, so the former does too. As $b_\cD(z)=1$, no
other path ends at $z$. Proposition~\ref{prop:EEequiv} finishes the proof.
\end{proof}

The hypotheses of Proposition~\ref{prop:moveA} and Theorem~\ref{thm:triangle} are avoidance statements
in graphs with at most one even vertex $g$, for a path $Q_c$ and a set $A$ that contains no vertex of
$N[c]$ and none of $N(g)$. In Proposition~\ref{prop:moveA}: $c=z$, $A=\{u\}$, $g=w$. In
Theorem~\ref{thm:triangle}: $c=w$, $A=\{u,z\}$, $g=z$. EE itself, through
Proposition~\ref{prop:EEequiv}, is the case $c=z$, $A=\{g\}=\{u\}$ in $K-uz$.

\begin{conjecture}[$L_1^*$]\label{conj:L1star}
Let $G$ have at most one vertex $g$ of positive even degree, let $c$ be odd, and let $A$ be a set of at
most two vertices with $A\cap(N[c]\cup N(g))=\emptyset$ ($g\in A$ allowed). Then some Lov\'asz
decomposition of $G$ has $Q_c\cap A=\emptyset$.
\end{conjecture}

For $|A|=1$ and no even vertex this is $L_1$~\cite[Theorem~6.3]{RD}. For $A=\{g\}$ it is EE. For
$A=\{a\}$ with $a\ne g$ it is $L_1$ in the presence of an even vertex. This fails when $g\sim a$ (on the
path $c\,g\,a$, say)~\cite[Remark~6.7]{RD}. The condition $a\notin N(g)$ is exactly what our
computations show to be needed: every one of the $6{,}423$ failures of $L_1$ with an even vertex at
order $9$ has $g\sim a$.

\begin{corollary}\label{cor:EEfromL1}
Conjecture~\ref{conj:L1star} implies Conjecture~\ref{conj:ee}.
\end{corollary}

\begin{proof}
Let $T=N(u)\cap N(z)$. If $|T|\le1$, use Theorem~\ref{thm:common}. If some $w\in N(u)$ lies outside
$N[z]$, apply Conjecture~\ref{conj:L1star} to $K-uz-uw$, $c=z$, $A=\{u\}$, and then
Proposition~\ref{prop:moveA}. Here $u\notin N[z]$, $u\notin N(w)$, and $w$ is the even vertex.
Otherwise take $w\in T$ and apply the conjecture to $K_1=K-\{uz,uw,wz\}$, $c=w$, $A=\{u,z\}$; then use
Theorem~\ref{thm:triangle}. Here $u,z\notin N_{K_1}[w]$, the even vertex is $z$ (or there is none),
and $u\notin N_{K_1}(z)$.
\end{proof}

\begin{remark}
Two stronger invariants fail. First, separation of $u$ and $z$ in $K-uz$, which Theorem~\ref{thm:common}
provides when $|T|\le1$: it fails for most configurations with $|T|\ge2$ (at order $7$: $99$ of $135$
with $|T|=2$, and all $65$ with $|T|\ge3$). Second, separation of $w$ from $\{u,z\}$ in $K_1$: it fails
in $267$ of $622$ triangles at order~$7$. So the avoidance in Conjecture~\ref{conj:L1star} is one-sided,
like $L_1$, and the separation theorem cannot supply it directly. The leaf trick that derives $L_1$ from
separation would make $c$ a second even vertex.
\end{remark}

\section{Computations}\label{sec:comp}

\emph{This section is computational. No proof above depends on it.} Scripts are in
\texttt{scripts/}, outputs in \texttt{scripts/out/}.

\subsection*{Oracle and validation}
\texttt{pee.c} decides whether $xy$ is terminal at $x$ in some Lov\'asz decomposition. It builds $Q_x$
first, as $x,y,\dots$ ending at an odd vertex, and then searches exactly for a decomposition of the rest
with the remaining end counts. It branches on the paths starting at an odd vertex of least remaining
degree, memoises failures, and prunes by two necessary conditions: every non-trivial component has an
end, and $d(v)\le2P-b(v)$ with $P$ the number of paths still to choose. Further modes decide avoidance
($Q_c\cap A=\emptyset$) and separation. Validation:
\begin{itemize}
\item against \texttt{brute.py}, an independent enumeration of all transition systems (at each
vertex a pairing of the incident edges, the unpaired edge being the end edge): all graphs of order
$\le7$ with at most one even vertex, $2{,}279$ queries, $0$ disagreements;
\item against \texttt{tools/route\_ext.c}, through Proposition~\ref{prop:EEequiv} (query \texttt{X}) and
Lemma~\ref{lem:leafequiv} (query \texttt{S}): order $9$, $47{,}688$ EE and $47{,}738$ away queries,
$0$ disagreements;
\item no failure on the $224$ connected all-odd graphs of order $8$ ($6{,}359$ queries), as
Theorem~\ref{thm:endedge} requires.
\end{itemize}

\subsection*{Exhaustive results}
All connected graphs with exactly one vertex of positive even degree (\texttt{nauty-geng -c}; such
graphs have odd order). ``Away'' counts the ordered edges $(x,y)$ with $y\ne z$ and $d(x)\ge3$.

\begin{center}
\begin{tabular}{rrrrrr}
\toprule
$n$ & graphs & EE tests & EE failures & away tests & away failures \\
\midrule
5 & 7 & 14 & 0 & 37 & 5 \\
7 & 130 & 401 & 0 & 1{,}910 & 54 \\
9 & 11{,}936 & 47{,}688 & 0 & 330{,}449 & 1{,}410 \\
11 & \multicolumn{5}{l}{EE only: see \texttt{scripts/out/ee11\_*.txt} and below} \\
\bottomrule
\end{tabular}
\end{center}

At order $11$ the EE run covers every connected graph with one even vertex. The final counts are in
\texttt{research/experiments.md} (E23); at the time of writing more than $30$ million $(z,u)$ pairs had
been tested, with no failure. A partial order-$11$ run on all edges (an initial segment of the
enumeration) found $56{,}739$ away failures and no EE failure.

\emph{The structure of the away failures.} Every failure at orders $5$, $7$, $9$ has $y\notin N(z)$ and
$\mathrm{dist}(x,z)\le2$, and in every one $xy$ is terminal at $y$. Every failure at orders $7$, $9$ and in the partial
order-$11$ run is unreachable in the sense of Section~\ref{sec:away}. At order $9$ the reachable away
edges number $270{,}875$, with no failure; the unreachable ones number $59{,}574$, of which $1{,}410$
fail.

\emph{$L_1^*$.} At order $9$ (all graphs with one even vertex): $L_1$ with an even vertex $g\not\sim a$ (and $a$ odd), $161{,}652$
tests, $0$ failures; with $g\sim a$, $6{,}423$ failures in $170{,}782$ tests. The case $|A|=2$ of
Conjecture~\ref{conj:L1star}: $203{,}658$ tests, $0$ failures. The hypothesis of
Theorem~\ref{thm:triangle} held for every one of the $88{,}356$ triples $(z,u,w)$ (not merely for some
$w$).

\subsection*{Random graphs}
\texttt{rand\_test.py 6000 101 5000000 11 13 15 17}: $6{,}000$ connected graphs of orders
$11$--$17$ with one even vertex ($G(n,p)$, $p\in[0.15,0.5]$, parities repaired by toggling edges),
node budget $5\cdot10^6$ per query.

\begin{center}
\begin{tabular}{lrrr}
\toprule
statement & tests & failures & undecided \\
\midrule
EE, $|N(u)\cap N(z)|\le1$ & $13{,}200$ & 0 & 120 \\
EE, $|N(u)\cap N(z)|\ge2$ & $12{,}762$ & 0 & 112 \\
away edges (sample of $12$ per graph) & $71{,}965$ & 99 & 457 \\
\quad of which $y\in N(z)$, or $\mathrm{dist}(x,z)\ge3$, or both ends fail & & 0 & \\
$L_1^*$, $A=\{a\}$, $a\ne g$ & $42{,}315$ & 0 & 276 \\
$L_1^*$, $A=\{g\}$ & $5{,}669$ & 0 & 27 \\
hypothesis of Theorem~\ref{thm:triangle} & $44{,}892$ & 0 & 383 \\
\bottomrule
\end{tabular}
\end{center}

\section{Formalisation}\label{sec:lean}

\texttt{lean/GallaiLean/PrescribedEnd.lean} (Lean~4, Mathlib, no \texttt{sorry}; \texttt{\#print
axioms} shows only \texttt{propext}, \texttt{Classical.choice}, \texttt{Quot.sound}; the full library
builds). A Lov\'asz decomposition is a \texttt{PathDecomp} with \texttt{ends v = degree v \% 2}, and
``$xy$ terminal at $x$'' is $\exists P\in D,\ (P.\mathtt{src}=x\lor P.\mathtt{tgt}=x)\land
s(x,y)\in P.\mathtt{edges}$.
\begin{itemize}
\item \texttt{sep\_atMostOneCommon}, \texttt{ee\_atMostOneCommon}: Theorem~\ref{thm:common};
\texttt{ee\_deg\_two}: Corollary~\ref{cor:deg2} (first statement).
\item \texttt{PathDecomp.append\_move\_exact}, \texttt{PathDecomp.liftLeaf},
\texttt{endEdge\_of\_leafSep}: the passage to the leaf graph and back (Lemmas~\ref{lem:leafequiv}
and~\ref{lem:merge}).
\item \texttt{pee\_adjEven}: Corollary~\ref{cor:adj}, case $k=1$, proved as in Lemma~\ref{lem:protect}
inside the leaf graph.
\item \texttt{ee\_of\_triangle}: Theorem~\ref{thm:triangle}.
\end{itemize}
Theorem~\ref{thm:path} for $k\ge2$, Proposition~\ref{prop:pendant}, Proposition~\ref{prop:moveA}
and Proposition~\ref{prop:counter} are proved here only. The library supplies Lov\'asz's theorem, the
separation theorem (\texttt{separation}), Fan's lemma with vertex sets (\texttt{fan\_single\_supp}),
the leaf merge (\texttt{leaf\_merge}) and Proposition~\ref{prop:EEequiv} (\texttt{ee\_of\_avoid}).

\section{Open problems}\label{sec:open}
\begin{enumerate}
\item Conjecture~\ref{conj:ee} when $u$ and $z$ have at least two common neighbours; by
Corollary~\ref{cor:EEfromL1} it suffices to prove Conjecture~\ref{conj:L1star}.
\item The case $|A|=1$, $A\ne\{g\}$, of Conjecture~\ref{conj:L1star} ($L_1$ with an even vertex not adjacent
to the avoided vertex). It is the hypothesis of Proposition~\ref{prop:moveA}.
\item Question~\ref{q:away}, and a characterisation of the unreachable edges that fail.
\item Several prescriptions at once (the star-terminal programme~\cite{STL}): every path through $z$
of the form $x\,z\,y\cdots$ with $x$ an end. Theorem~\ref{thm:common} prescribes one end edge at $z$.
Lemma~\ref{lem:protect} is the tool for keeping it while others are added: a prescription survives
every Fan transfer whose target is not adjacent to its end, or is its other end.
\end{enumerate}

\begin{thebibliography}{9}
\bibitem{Fan05} G.~Fan, Path decompositions and Gallai's conjecture, J.~Combin.\ Theory Ser.~B 93 (2005)
117--125.
\bibitem{L68} L.~Lov\'asz, On covering of graphs, in: Theory of Graphs (Proc.\ Colloq.\ Tihany, 1966),
Academic Press, New York, 1968, 231--236.
\bibitem{MS} J.~Williams, Separating many pairs in Lov\'asz decompositions: classes, the leaf merge, and
the Case~I step, manuscript (\texttt{manuscripts/multi-sep}), 2026.
\bibitem{RD} J.~Williams, Route descent: transversal realisability, minimum $T$-joins, and Lov\'asz
decompositions that separate two vertices, manuscript (\texttt{manuscripts/route-descent}), 2026.
\bibitem{STL} J.~Williams, Star-terminal Lov\'asz decompositions and the spine tests, manuscript
(\texttt{manuscripts/stl}), 2026.
\end{thebibliography}
\end{document}
